\subsection{Idris2}
\seclabel{idris2}
At its heart, the type checker implements: dynamic
pattern
unification~\cite{miller:unification,reed:unification,gundry:thesis},
which instantiates implicit arguments; and conversion checking
whether two terms evaluate to the same reduct.
Both components require an evaluator.  \idris{} uses
normalisation by evaluation~\cite{berger:nbe} with a \emph{syntactic representation}
(terms) and a
\emph{semantic representation} (values in weak head normal form). The
static evaluator is call-by-name and produces a weak head normal form from a term, and
\idris{} implements a quotation mechanism which reconstructs a term from a semantic
representation of a weak head normal form.

Profiling the \idris{} executable reveals that most performance
bottlenecks encountered in developing \Frexlib{} are caused by
the evaluator.
%
We have experimented with alternative implementations of the evaluator
that compile terms using Scheme's backend and a glued
representation of values~\cite{DBLP:journals/mscs/CoquandD97,Chapman2005epigram} rather
than interpreting terms directly.
%
These implementations give
modest performance
gains, but in the end the most effective way is to avoid evaluation in the first place! We have therefore also experimented with the following ways
to avoid evaluating terms: preserving subterm sharing, choosing appropriate
data representation in unification, and taking advantage of the typical
structure of unification and conversion problems.

\subsubsection*{Preserving Sharing}

Instantiating implicit arguments in dependently typed programs often
leads to significant \emph{sharing} of subterms.
For example,
\IdrisData{[}\IdrisData{True}\IdrisData{,} \IdrisData{False}\IdrisData{]}
\IdrisKeyword{:} \IdrisType{Vect} \IdrisData{2} \IdrisType{Bool}
elaborates to the following term:
\IdrisKeyword{(}\IdrisData{::}\IdrisKeyword{)} \IdrisKeyword{(}\IdrisData{S} \IdrisData{Z}\IdrisKeyword{)} \IdrisType{Bool} \IdrisData{True} \IdrisKeyword{(}\IdrisKeyword{(}\IdrisData{::}\IdrisKeyword{)} \IdrisData{Z} \IdrisType{Bool} \IdrisData{False} \IdrisKeyword{(}\IdrisData{Nil} \IdrisType{Bool}\IdrisKeyword{)}\IdrisKeyword{)},
sharing the subexpressions \IdrisData{Z} and \IdrisType{Bool}. As the vector
gets longer, sharing increases.
Following~\citet{kovacs:smalltt}, we preserve sharing by introducing a
metavariable for \emph{every} implicit argument, inlining only when it is
guaranteed that the definition cannot break sharing. Consequently,
we inline a metavariable whose definition is itself a metavariable applied
to local variables. Otherwise, we do not substitute metavariable solutions
into terms at all until they are required for unification or display
purposes.

\subsubsection*{Unification}

Unification operates on \emph{values}, not \emph{terms}, but sometimes \idris{}
needs to postpone a unification problem if it is blocked due to an unsolved
metavariable. When the metavariable is solved, \idris{} needs to re-evaluate the
terms being unified. Previously, \idris{} stored postponed problems as a pair
of (syntactic) terms in an environment, re-evaluated once the
blocking metavariable is solved.
However, \Frexlib{} produces some large postponed problems, for which quotation
to syntax is expensive.  Now, in addition to the evaluator and quotation, we
have introduced a \emph{continue} operation, which re-evaluates the
metavariable at the head of a blocked value, and avoids unnecessary quotation.

\subsubsection*{Conversion Checking}

Types in \Frexlib{} can be large, and sometimes a unification problem that arises
while type checking \Frexlib{} is postponed due to an unsolved
metavariable which blocks evaluation. In this case, we might have a
unification problem of the form
\IdrisFunction{f} \IdrisBound{x1} \IdrisKeyword{...} \IdrisBound{xn}
\IdrisKeyword{=?=}
\IdrisFunction{f} \IdrisBound{y1} \IdrisKeyword{...} \IdrisBound{yn}
where the \IdrisBound{xi}, \IdrisBound{yi} etc may be very large subterms,
and the terms unify if they are convertible.
If most corresponding terms are equal after evaluation, but one differs, it may
take a long time to find the differing subterm which blocks unification,
especially since checking the convertibility of subterms involves evaluation.
Fortunately, terms in blocked unification problems
tend to differ at the heads rather than at deeply nested subterms. Therefore,
we always check the heads of the values of corresponding \IdrisBound{xi} and
\IdrisBound{yi} first, postponing the unification
problem if any are unequal. This heuristic significantly improves performance,
preventing a lot of unnecessary evaluation.

\subsubsection*{Influence on Language Design and Ecosystem}

The development of \Frexlib{} has led to the implementation of a number of
desirable language features in \idris{}.
%
Many of these have been minor changes to the treatment of implicit
arguments and \texttt{parameters} blocks.
%
More significantly, \Frexlib{} makes extensive use of \texttt{auto}
implicit arguments, which are solved by a search procedure which uses
constructors and functions marked as search hints.  To assist
the development of \Frexlib{} and improve the readability of its code we have added the
ability to mark \emph{local} functions as search hints, allowing
us to restrict the scope of the hints and so avoid excessive
growth of the search space.
%
\Frexlib{} is now part of the \idris{} test suite, ensuring that it will
remain consistent with any updates to \idris{}.
